\scanpage{89}
et que $\inserted{G\mapsto}S\Pi M_G$ est un foncteur en \deleted{$\unsure{M}_G$} $G$ pour les morphismes quelconques, et que (124) est fonctoriel pour les iso. Pour un morphisme quelconque $G\xrightarrow{f}G'$, correspondant à une partie $C\subset A_G$ (telle que $G\unsure{\rightharpoondown|{-}}G'$ soit isom, en tant que quotient de $G$, à $G\unsure{\rightharpoondown|{-}}G/C$), on suppose que l'on a un diagramme commutatif
\[
\tag{125}
\diagram{\begin{tikzcd}[ampersand replacement=\&, column sep=large,row sep=large]
S\Pi M_G \arrow[r,"S\Pi(f)"] \arrow[d,"\varphi_G"] \& S\Pi M_{G'} \arrow[d,"\varphi_{G'}"] \\
S_{\deleted{\unsure{A}}C}M_G \arrow[r,"\gamma_f"] \& \deleted{\unsure{S}}\Pi M_{G'}
\end{tikzcd}}
\]
i.e. que le composé
\[
\diagram{\begin{tikzcd}[ampersand replacement=\&, column sep=large]
S\Pi M_G \arrow[r,"S\Pi(f)"] \& S\Pi M_{G'} \arrow[r,phantom,"{\overset{\varphi_{G'}}{\unsure{\rightharpoondown|{-}}}}" description] \& \Pi M_{G'}
\end{tikzcd}}
\]
se factorise par l'extension $S_{\deleted{\unsure{A}}C}M_G$ de $\Pi M_G$ par $\mathbb Z^C$, déduite de l'extension $S\Pi M_G$ de $\Pi M_G$ par $\mathbb Z^{A(G)}$ grâce à l'hom canonique
\[
\diagram{\begin{tikzcd}[ampersand replacement=\&]
\mathbb Z^{A(G)} \arrow[r] \& \mathbb Z^C.
\end{tikzcd}}
\]
On a automatique alors la relation de transitivité pour un composé.
\margin{88}
